% Template laporan ini dibuat dengan harapan untuk standardisasi laporan praktikum pada lingkungan DTEDI SV UGM, agar semua laporan praktikum konsisten secara format dan hasil cetak laporan.
% Saran dan perbaikan sangat diperlukan agar template ini dapat digunakan oleh semua orang. bila ada saran, silahkan hubungi author melalui linkedin : Alim Satria
% Dibuat oleh : Alim Satria Fi'i Wijaya Kusuma

\documentclass[12 pt]{article}
\usepackage[a4paper, total={6in, 9in}]{geometry}
\usepackage[utf8]{inputenc}
\usepackage{color}
\usepackage[bahasa]{babel}
\usepackage{float}
\usepackage{graphicx}
\usepackage{setspace}
\usepackage{indentfirst}
\usepackage{subfiles}
\usepackage{mathtools}
\usepackage{tikz}
\usepackage{pgfplots}

\renewcommand \thesection{\Roman{section}.}
\renewcommand \thesubsection{\arabic{section}.\arabic{subsection}.}
\renewcommand \thesubsubsection{\arabic{section}.\arabic{subsection}.\arabic{subsubsection}.}
\renewcommand \figurename{Gambar}
\renewcommand \thefigure {\arabic{section}.\arabic{figure} } %kustom nomor gambar dengan format (section.subsection.urutan_figure)
\renewcommand \tablename{Tabel}
\renewcommand \thetable {\arabic{section}.\arabic{subsection}.\arabic{table} } %kustom nomor tabel dengan format (section.subsection.urutan_tabel)
\renewcommand \theequation{\arabic{section}.\arabic{equation}} %kustom pemomoran persamaan matematika (section.urutan_persamaan)
\setlength{\tabcolsep}{10pt} %kustom spacing coloumn tabel (default = 6pt)
\renewcommand{\arraystretch}{1.5} %kustom spacing row tabel (default = 1)
\newcommand{\ffrac}[2]{\ensuremath{\frac{\displaystyle #1}{\displaystyle #2}}}

\title{\large \textbf{LAPORAN PRAKTIKUM TEKNIK KENDALI LANJUT} \\ \textbf{MODUL 1 DAN LATIHAN SOAL}\linebreak}
\author{}
\date{}
\onehalfspacing
\pgfplotsset{compat = newest}

\begin{document}

\clearpage
\maketitle
\thispagestyle{empty} %menghapus penomoran pada cover

\begin{center}
\includegraphics[width=5cm,height=5cm]{cover_laporan/logo_ugm.png}
\end{center}

\vspace{1 cm}

\begin{center}
\begin{tabular}{lcl}
Nama & : & Alim Satria Fi'i Wijaya Kusuma \\
NIM & : & 21/483503/SV/20304\\
Kelas & : & IKPAB\\
Tanggal & : & 29 Agustus 2022 dan 5 September 2022 \\
Dosen Pengampu & : & Fahmizal, S.T., M.Sc.
\end{tabular}
\newline
\newline
\newline
\newline
\newline
\large{\textbf{PROGRAM STUDI TEKNOLOGI REKAYASA INSTRUMENTASI DAN KONTROL}} \\
\textbf{SEKOLAH VOKASI UNIVERSITAS GADJAH MADA} \\
\textbf{2022}
\end{center}

\pagebreak

\subfile{sections/isi.tex}

\end{document}